\input{../tex-inputs/latex-leseansicht-vorspann}

\section[Arthur Schnitzler und Olga und Elisabeth Gussmann an Gustav Schwarzkopf, 14. 8. 1901]{L04436 Arthur Schnitzler und Olga und Elisabeth Gussmann an Gustav
               Schwarzkopf, 14. 8. 1901}
\nopagebreak\mylabel{L04436v}
\rehead{ }\normalsize\beginnumbering\briefempfaengerindex{Schwarzkopf, Gustav@\textsc{Schwarzkopf, Gustav}!zzzSteinrück, Elisabeth@\emph{von Elisabeth Steinrück}!1901-08-141@{14. 8. 1901}|(be}\briefempfaengerindex{Schwarzkopf, Gustav@\textsc{Schwarzkopf, Gustav}!zzzSchnitzler, Olga@\emph{von Olga Schnitzler}!1901-08-141@{14. 8. 1901}|(be}\briefempfaengerindex{Schwarzkopf, Gustav@\textsc{Schwarzkopf, Gustav}!zzzSchnitzler, Arthur@\emph{von Arthur Schnitzler}!1901-08-141@{14. 8. 1901}|(be}
\toendnotes[C]{\smallbreak\pagebreak[2]}
\correspDesc{Versand  durch Arthur Schnitzler, Olga Gussmann, Elisabeth Gussmann am 14. 8. 1901 in Ritten
\newline{}Übermittlung  am 14. 8. 1901 in Bozen
\newline{}Zustellung  am 15. 8. 1901 in Wien
\newline{}Erhalt  durch Gustav Schwarzkopf am 15. 8. 1901 in Wien}\toendnotes[C]{\smallbreak}
\Standort{DLA, A:Schnitzler, HS.1985.1.1897.}
\physDesc{Bildpostkarte, 104 Zeichen
\newline{}Handschrift Arthur Schnitzler: Bleistift, deutsche Kurrent
\newline{}Handschrift Olga Schnitzler: Bleistift, deutsche Kurrent
\newline{}Handschrift Elisabeth Steinrück: Bleistift, deutsche Kurrent
\newline{}Versand: 1) Stempel: »\nobreak{}14. 8. 0\textcolor{gray}{1}, \oindex{Bozen@\textbf{Bozen}|pwk}\textcolor{gray}{Bozen}, \textcolor{gray}{2}N\nobreak{}«.   2) Stempel: »\nobreak{}\oindex{I., Innere Stadt@\textbf{I., Innere Stadt}|pwk}Wien 1/1 1, 15. 8. 01, 11½–1V., Bestellt\nobreak{}«. }\pstart{}{\pb}Herrn Gustav  Schwarzkopf\pend{}\pstart{}Wien\oindex{Wien@\textbf{Wien}|pw}\pend{}\pstart{}I Tiefer Graben 23\oindex{Wien@\textbf{Wien}!I., Innere Stadt@\textbf{I., Innere Stadt}!Tiefer Graben 23@\textbf{Tiefer Graben 23}|pw}.\pend{}{\bigskip}
\pstart
           \noindent{}\centering{}{\pb}\textcolor{gray}{\textbf{Runkelstein bei Bozen}}\oindex{Castel Roncolo@\textbf{Castel Roncolo}|pw}\pend
           \vspace{1em}
\pstart
           \raggedleft{}{\pb}14/8 901\pend
           \vspace{0.5em}
\pstart
           Herzliche Grüße\pend
           \pstart Ihr \spacefill\mbox{Arthur}\pend{}\selectlanguage{ngerman}\vspace{1em}
\pstart
           {[}hs. O. Schnitzler:{]} Herzl.{\\[\baselineskip]}\spacefill\mbox{Olga}\pend
           \leftskip=0em{}\selectlanguage{ngerman}\vspace{1em}
\pstart
           {[}hs. Steinrück:{]} Schöne Grüße{\\[\baselineskip]}\spacefill\mbox{L.}\pend
           \leftskip=0em{}\selectlanguage{ngerman}\endnumbering\briefempfaengerindex{Schwarzkopf, Gustav@\textsc{Schwarzkopf, Gustav}!zzzSteinrück, Elisabeth@\emph{von Elisabeth Steinrück}!1901-08-141@{14. 8. 1901}|)be}\briefempfaengerindex{Schwarzkopf, Gustav@\textsc{Schwarzkopf, Gustav}!zzzSchnitzler, Olga@\emph{von Olga Schnitzler}!1901-08-141@{14. 8. 1901}|)be}\briefempfaengerindex{Schwarzkopf, Gustav@\textsc{Schwarzkopf, Gustav}!zzzSchnitzler, Arthur@\emph{von Arthur Schnitzler}!1901-08-141@{14. 8. 1901}|)be}\mylabel{L04436h}
\begin{anhang}
\end{anhang}\newcommand{\dateiname}{L04436}\newcommand{\titel}{Arthur Schnitzler und Olga und Elisabeth Gussmann an Gustav Schwarzkopf, 14. 8. 1901}\newcommand{\editorInnen}{Herausgegeben von Jahnke, SelmaMüller, Martin Anton}\input{../tex-inputs/latex-leseansicht-abspann}
